% ===== 7 结论 =====
\section{结论}
\label{sec:r-conclusion}

本文面向自然堆积混凝土骨料的智能筛分任务，构建了一条从原始图像到标准筛分语义输出的可解释处理链。方法首先利用 ImageGrains 2.0 / Cellpose-SAM 的预训练实例感知能力恢复独立颗粒，再通过本次成像尺度标定获得毫米级几何量；在此基础上，用可校准的筛分等效粒径处理二维投影与筛孔语义差异，用 $w=d^{\gamma}$ 处理颗粒数量统计与机械筛分质量统计之间的口径差异，最终输出标准粒级质量占比和 $D_{10}/D_{50}/D_{90}$。同一实例几何还被复用于投影形貌和异常候选，从而形成可追溯的场景级报告。

现有三张真实自然堆积样本已经验证端到端功能闭环，并显示数量型与质量代理型级配存在显著差异：数量 $D_{50}$ 为 5.2--7.6~mm，而当前 $d^3$ 质量代理下为 14.1--27.5~mm。这一结果支持“标准筛分不能直接由颗粒个数分布替代”的核心建模判断；27 项竞赛层合成测试和约 27~s/张的历史算法时序进一步支持软件一致性与计算可行性。但当前仍缺与演示图配对的机械筛分逐筛称量真值，因此本文没有把未校准视觉输出表述为筛分准确率。

下一阶段最关键的工作是完成同批图像--机械筛分配对实验：用校准批次拟合并冻结 $\boldsymbol{\theta},\gamma$，再在独立批次上报告 $D_p$ 误差、逐筛级质量误差及核心消融；同时补充人工实例 mask，将最终误差进一步分解为视觉感知与筛分语义映射两部分。若这些独立验证表明低维映射能够稳定逼近标准筛分，则“预训练实例感知 + 显式物理/统计校准”将提供一条兼顾现场部署、可解释性与可扩展性的视觉筛分路线；若残差主要来自遮挡和第三维缺失，则后续升级应转向多视角/深度，而不是继续增加二维规则复杂度。
